\section{System Design}\label{sec:design}
Figure~\ref{fig:arch} shows the control loop. A health monitor samples CPU
temperature, throttle flags (\texttt{vcgencmd get\_throttled} on Pi, with a
sysfs/\texttt{psutil} fallback), memory, swap, CPU utilisation, frame rate and
camera health once per second. A deterministic decision engine maps each sample
to a rung of the ladder; a model manager realises the change.

\begin{figure}[t]
\centering
\begin{tikzpicture}[node distance=5mm and 6mm, every node/.style={font=\scriptsize},
  box/.style={draw,rounded corners,minimum height=6mm,minimum width=14mm,align=center},
  >=Stealth]
\node[box] (cam) {Camera /\\source};
\node[box,right=of cam] (pipe) {Frame loop\\+ scheduler};
\node[box,right=of pipe] (det) {Detector\\(active rung)};
\node[box,below=of pipe] (hm) {Health\\monitor};
\node[box,below=of det] (eng) {Decision\\engine};
\node[box,below=of hm] (log) {Audit log\\(SQLite)};
\node[box,below=of eng] (mm) {Model manager\\(preload, swap)};
\node[box,left=of log] (wd) {Watchdog\\supervisor};
\draw[->] (cam)--(pipe); \draw[->] (pipe)--(det);
\draw[->] (hm)--(eng); \draw[->] (eng)--(mm); \draw[->] (eng)--(log);
\draw[->] (mm.east) -- ++(4mm,0) |- (det.east);
\draw[->] (pipe)--(hm);
\draw[<->,dashed] (wd) -- (pipe);
\end{tikzpicture}
\caption{\shade{} architecture. The decision engine never touches models; the
model manager swaps a preloaded backend between frames. A separate supervisor
process watches a heartbeat from the frame loop.}\label{fig:arch}
\end{figure}

\subsection{Degradation ladder}
A rung is a tuple (model, input size, detection interval $N$, threads,
backend). Rungs are ordered best to worst; rungs whose model file is absent are
skipped. Moving between rungs that share a model is free; otherwise the target
model is loaded on a low-priority thread (with a RAM-headroom check and a warm-up
inference) and published for an atomic reference swap at the next frame
boundary. The latest request wins; a failed load marks the rung unavailable.

\subsection{Decision policy}
The engine reads a pressure level $\ell\in\{0,1,2,3\}$ from temperature bands
(warm/degrade/critical), firmware flags, memory/swap and optional frame-rate and
camera signals, plus an anticipation term when the temperature slope over a
window exceeds a threshold while warm. $\ell$ maps to a minimum rung. Moves
\emph{down} need \texttt{confirm\_samples} consecutive samples and a dwell
time, except at critical pressure; moves \emph{up} are one rung at a time, only
after a calm hold below an upgrade gate with no throttle flags and a cooldown.
All thresholds are configuration (Table~\ref{tab:params} in the appendix).
Under-voltage is reported as a power fault rather than answered with a smaller
model.

\subsection{Supervision and audit}
The frame loop writes a heartbeat (also while idle). A supervisor restarts the
engine on exit or silence, with exponential backoff, starting on the emergency
rung or the last rung that ran stably (last-known-good), and gives up after
repeated unstable restarts so the init system can take over. Every decision,
swap, camera reconnect and restart is stored with its reasons and signal
values, so a deployment can be reviewed after the fact.

\subsection{Optional: thermal-forecast advisor}
As an experimental alternative to fixed thresholds, an advisor fits
$\dot T=(T_{\mathrm{amb}}+k\,u-T)/\tau$ online ($u$ is CPU utilisation) by
Bayesian recursive least squares with a physical prior and forgetting, calibrates
per-rung utilisation from measurements, and selects the best rung whose upper
confidence bound of the temperature $H$ seconds ahead stays below a target. Samples
taken while the firmware throttles are excluded from the fit. We evaluate it
in Sec.~\ref{sec:sim}.
